% GENERATED FILE. Edit canonical lessons, then run npm run generate:books.
% canonical-source-hash: fnv1a64:6d88308bcca0384b
% canonical-lessons: JA-C176-shouyu, JA-C176-miso, JA-C176-waru, JA-C176-musu, JA-C176-itameru

\chapter{Frying, Steaming and Seasoning}
\label{ch:ja-frying-steaming-and-seasoning}
\hlchaptermodality{176}

\begin{chapteropening}
\textbf{By the end of this chapter:} \emph{I can talk about seasonings and ways of cooking.}

Complete the last lesson of chapter 176: I can talk about seasonings and ways of cooking.
\end{chapteropening}

\section[\texorpdfstring{\ja{しょうゆ} (shōyu)}{しょうゆ (shōyu)}]{\ja{しょうゆ} (shōyu) — soy sauce}

\label{lesson:JA-C176-shouyu}



\begin{quote}
Before the new one: say the Japanese for rice bowl, then the Japanese for oil.
\end{quote}

\subsection*{You'll want to know: \ja{しょうゆ}}
\textbf{\ja{しょうゆ}} — \emph{shōyu} — \textquotedblleft{}soy sauce\textquotedblright{}.

\begin{quote}
\textbf{\ja{しょうゆを} \ja{すこし} \ja{つかいます。}} — \emph{shōyu o sukoshi tsukaimasu} — \textquotedblleft{}I use a little soy sauce.\textquotedblright{}
\end{quote}

\begin{cousinweb}
\textbf{\ja{しょうゆ}} (\emph{shōyu}) is Sino-Japanese (kango): two Chinese characters read with their on'yomi, readings borrowed from Chinese. They are \emph{shō}, \textquotedblleft{}fermented paste\textquotedblright{} (Middle Chinese \emph{tsjangH}) and \emph{yu}, \textquotedblleft{}oil\textquotedblright{} (\emph{yuw}) (KANJIDIC2; Middle Chinese in Baxter's transcription).
\end{cousinweb}

\begin{grammarlens}[title={What you've built}]
One more word for everyday life at A2.
\end{grammarlens}

\subsection*{Your turn}
\begin{itemize}
\raggedright
  \item \emph{Say it:} \emph{shōyu}
  \item \emph{Say it:} \emph{shōyu}, once more
  \item \emph{Say it:} \emph{shōyu o sukoshi tsukaimasu} — I use a little soy sauce.
  \item \emph{Recall:} say the Japanese for rice bowl, then the Japanese for oil, then say \emph{shōyu} again
\end{itemize}

\begin{tcolorbox}[breakable,colback=teal!4,colframe=teal!35!black,title={Before you move on}]
What does \ja{しょうゆ} mean? (\textquotedblleft{}Soy sauce\textquotedblright{}.) Say it once more.
\end{tcolorbox}
\section[\texorpdfstring{\ja{みそ} (miso)}{みそ (miso)}]{\ja{みそ} (miso) — soybean paste; miso}

\label{lesson:JA-C176-miso}



\begin{quote}
Before the new one: say the Japanese for oil, then the Japanese for soy sauce.
\end{quote}

\subsection*{You'll want to know: \ja{みそ}}
\textbf{\ja{みそ}} — \emph{miso} — \textquotedblleft{}soybean paste; miso\textquotedblright{}.

\begin{quote}
\textbf{\ja{みそと} \ja{とうふを} \ja{かいます。}} — \emph{miso to tōfu o kaimasu} — \textquotedblleft{}I buy miso and tofu.\textquotedblright{}
\end{quote}

\begin{cousinweb}
\textbf{\ja{みそ}} (\emph{miso}) is written with two characters read \emph{mi}, \textquotedblleft{}taste\textquotedblright{} (Middle Chinese \emph{mj+jH}), and \emph{so}; JMdict marks at least one of its senses as usually written in kana, and the word's older origin is not settled (JMdict, KANJIDIC2).
\end{cousinweb}

\begin{grammarlens}[title={What you've built}]
One more word for everyday life at A2.
\end{grammarlens}

\subsection*{Your turn}
\begin{itemize}
\raggedright
  \item \emph{Say it:} \emph{miso}
  \item \emph{Say it:} \emph{miso}, once more
  \item \emph{Say it:} \emph{miso to tōfu o kaimasu} — I buy miso and tofu.
  \item \emph{Recall:} say the Japanese for oil, then the Japanese for soy sauce, then say \emph{miso} again
\end{itemize}

\begin{tcolorbox}[breakable,colback=teal!4,colframe=teal!35!black,title={Before you move on}]
What does \ja{みそ} mean? (\textquotedblleft{}Soybean paste; miso\textquotedblright{}.) Say it once more.
\end{tcolorbox}
\section[\texorpdfstring{\ja{わる} (waru)}{わる (waru)}]{\ja{わる} (waru) — to crack; to split, to divide}

\label{lesson:JA-C176-waru}



\begin{quote}
Before the new one: say the Japanese for soy sauce, then the Japanese for soybean paste.
\end{quote}

\subsection*{You'll want to know: \ja{わる}}
\textbf{\ja{わる}} — \emph{waru} — \textquotedblleft{}to crack; to split, to divide\textquotedblright{}.

\begin{quote}
\textbf{\ja{たまごを} \ja{ふたつ} \ja{わります。}} — \emph{tamago o futatsu warimasu} — \textquotedblleft{}I crack two eggs.\textquotedblright{}
\end{quote}

Polite \textbf{\ja{わります}}, \emph{warimasu}; a godan verb.

\begin{cousinweb}
\textbf{\ja{わる}} (\emph{waru}) is a native Japanese word (wago), inherited rather than borrowed from Chinese; its kanji was attached later as a kun'yomi (JMdict, KANJIDIC2).
\end{cousinweb}

\begin{grammarlens}[title={What you've built}]
One more word for everyday life at A2.
\end{grammarlens}

\subsection*{Your turn}
\begin{itemize}
\raggedright
  \item \emph{Say it:} \emph{waru}
  \item \emph{Say it:} \emph{waru}, once more
  \item \emph{Say it:} \emph{tamago o futatsu warimasu} — I crack two eggs.
  \item \emph{Recall:} say the Japanese for soy sauce, then the Japanese for soybean paste, then say \emph{waru} again
\end{itemize}

\begin{tcolorbox}[breakable,colback=teal!4,colframe=teal!35!black,title={Before you move on}]
What does \ja{わる} mean? (\textquotedblleft{}To crack; to split, to divide\textquotedblright{}.) Say it once more.
\end{tcolorbox}
\section[\texorpdfstring{\ja{むす} (musu)}{むす (musu)}]{\ja{むす} (musu) — to steam}

\label{lesson:JA-C176-musu}



\begin{quote}
Before the new one: say the Japanese for soybean paste, then the Japanese for to crack.
\end{quote}

\subsection*{You'll want to know: \ja{むす}}
\textbf{\ja{むす}} — \emph{musu} — \textquotedblleft{}to steam\textquotedblright{}.

\begin{quote}
\textbf{\ja{ばんごはんに} \ja{やさいを} \ja{むします。}} — \emph{bangohan ni yasai o mushimasu} — \textquotedblleft{}I steam vegetables for dinner.\textquotedblright{}
\end{quote}

Polite \textbf{\ja{むします}}, \emph{mushimasu}; a godan verb.

\begin{cousinweb}
\textbf{\ja{むす}} (\emph{musu}) is a native Japanese word (wago), inherited rather than borrowed from Chinese; its kanji was attached later as a kun'yomi (JMdict, KANJIDIC2).
\end{cousinweb}

\begin{grammarlens}[title={What you've built}]
One more word for everyday life at A2.
\end{grammarlens}

\subsection*{Your turn}
\begin{itemize}
\raggedright
  \item \emph{Say it:} \emph{musu}
  \item \emph{Say it:} \emph{musu}, once more
  \item \emph{Say it:} \emph{bangohan ni yasai o mushimasu} — I steam vegetables for dinner.
  \item \emph{Recall:} say the Japanese for soybean paste, then the Japanese for to crack, then say \emph{musu} again
\end{itemize}

\begin{tcolorbox}[breakable,colback=teal!4,colframe=teal!35!black,title={Before you move on}]
What does \ja{むす} mean? (\textquotedblleft{}To steam\textquotedblright{}.) Say it once more.
\end{tcolorbox}
\section[\texorpdfstring{\ja{いためる} (itameru)}{いためる (itameru)}]{\ja{いためる} (itameru) — to stir-fry; to fry}

\label{lesson:JA-C176-itameru}



\begin{quote}
Before the new one: say the Japanese for to crack, then the Japanese for to steam.
\end{quote}

\subsection*{You'll want to know: \ja{いためる}}
\textbf{\ja{いためる}} — \emph{itameru} — \textquotedblleft{}to stir-fry; to fry\textquotedblright{}.

\begin{quote}
\textbf{\ja{あぶらで} \ja{やさいを} \ja{いためます。}} — \emph{abura de yasai o itamemasu} — \textquotedblleft{}I stir-fry vegetables in oil.\textquotedblright{}
\end{quote}

Polite \textbf{\ja{いためます}}, \emph{itamemasu}; an ichidan verb.

\begin{cousinweb}
\textbf{\ja{いためる}} (\emph{itameru}) is a native Japanese word (wago), inherited rather than borrowed from Chinese; its kanji was attached later as a kun'yomi (JMdict, KANJIDIC2).
\end{cousinweb}

\begin{grammarlens}[title={What you've built}]
One more word for everyday life at A2.
\end{grammarlens}

\subsection*{Your turn}
\begin{itemize}
\raggedright
  \item \emph{Say it:} \emph{itameru}
  \item \emph{Say it:} \emph{itameru}, once more
  \item \emph{Say it:} \emph{abura de yasai o itamemasu} — I stir-fry vegetables in oil.
  \item \emph{Recall:} say the Japanese for to crack, then the Japanese for to steam, then say \emph{itameru} again
\end{itemize}

\begin{tcolorbox}[breakable,colback=teal!4,colframe=teal!35!black,title={Before you move on}]
What does \ja{いためる} mean? (\textquotedblleft{}To stir-fry; to fry\textquotedblright{}.) Say it once more.
\end{tcolorbox}
